\documentclass[11pt,a4paper]{article}

\usepackage{fontspec}
\setmainfont{DejaVu Serif}[
  BoldFont={DejaVu Serif Bold},
  ItalicFont={DejaVu Serif},
  ItalicFeatures={FakeSlant=0.18},
  BoldItalicFont={DejaVu Serif Bold},
  BoldItalicFeatures={FakeSlant=0.18}
]
\setsansfont{DejaVu Sans}
\usepackage{polyglossia}
\setmainlanguage{russian}
\setotherlanguage{english}
\newfontfamily\cyrillicfonttt{DejaVu Sans Mono}
\usepackage[a4paper,margin=21mm]{geometry}
\usepackage{microtype}
\usepackage{setspace}
\setstretch{1.04}
\usepackage{amsmath,amssymb,amsthm,mathtools,bm,mathrsfs}
\usepackage{booktabs,longtable,array,tabularx}
\usepackage{enumitem}
\usepackage{ragged2e}
\usepackage{tocloft}
\setlength{\cftsecnumwidth}{2.5em}
\setlength{\cftsubsecnumwidth}{3.4em}
\usepackage{xcolor}
\usepackage{tikz}
\usetikzlibrary{arrows.meta,positioning,calc,fit}
\usepackage[unicode,hidelinks]{hyperref}
\usepackage[nameinlink,noabbrev]{cleveref}
\emergencystretch=3em
\allowdisplaybreaks

\hypersetup{
  pdftitle={Functional Interface for Geometry of Observation: GO-QS Typed Variational Protocol},
  pdfauthor={Stassis Research Program},
  pdfsubject={P0 reference implementation of the GO-QS observation interface},
  pdfkeywords={geometry of observation, GO-QS, inverse problems, typed quantities, variational estimator}
}

\definecolor{obsblue}{RGB}{37,83,138}
\definecolor{obsred}{RGB}{150,45,45}
\definecolor{obsgreen}{RGB}{30,115,78}
\definecolor{obsgray}{RGB}{75,75,75}
\definecolor{lightblue}{RGB}{239,246,253}
\definecolor{lightred}{RGB}{253,242,242}
\definecolor{lightgreen}{RGB}{239,249,244}

\theoremstyle{definition}
\newtheorem{definition}{Определение}[section]
\newtheorem{model}[definition]{Модель}
\newtheorem{hypothesis}[definition]{Гипотеза}
\theoremstyle{plain}
\newtheorem{theorem}[definition]{Теорема}
\newtheorem{proposition}[definition]{Предложение}
\newtheorem{lemma}[definition]{Лемма}
\theoremstyle{remark}
\newtheorem{remark}[definition]{Замечание}
\newtheorem{warning}[definition]{Ограничение}

\newcommand{\R}{\mathbb{R}}
\newcommand{\N}{\mathbb{N}}
\newcommand{\B}{\mathcal{B}}
\newcommand{\A}{\mathcal{A}}
\newcommand{\Qreg}{\mathcal{Q}}
\newcommand{\Dprof}{\mathscr{D}}
\newcommand{\dd}{\mathop{}\!d}
\newcommand{\norm}[1]{\left\lVert #1\right\rVert}
\newcommand{\abs}[1]{\left\lvert #1\right\rvert}
\newcommand{\set}[1]{\left\{#1\right\}}
\newcommand{\given}{\,\middle|\,}
\newcommand{\Type}{\operatorname{Type}}
\newcommand{\Cod}{\operatorname{Cod}}
\newcommand{\Argmin}{\operatorname*{arg\,min}}
\newcommand{\dist}{\operatorname{dist}}
\newcommand{\rank}{\operatorname{rank}}
\newcommand{\rad}{\operatorname{rad}}
\newcommand{\id}{\operatorname{id}}
\newcommand{\status}[1]{\textsf{\footnotesize #1}}

\newcolumntype{L}[1]{>{\raggedright\arraybackslash}p{#1}}
\newcolumntype{Y}{>{\raggedright\arraybackslash}X}

\title{\bfseries Функциональный интерфейс геометрии наблюдения\\
\large Типизированный вариационный протокол GO-QS}
\author{Stassis Research Program}
\date{Строгая эталонная редакция v0.2 (P0) --- 2026}

\begin{document}
\setlength{\RaggedRightParindent}{1.2em}
\RaggedRight
\hyphenation{
ис-поль-зо-вать
един-ствен-ность
при-ви-ле-ги-ро-ван-ность
про-стран-ствен-ной
пре-об-ра-зо-ва-нии
без-раз-мер-на-я
су-ще-ство-ва-ни-е
не-един-ствен-ность
ин-тер-пре-та-ци-я
ре-флек-сив-но-е
по-сле-до-ва-тель-но
ко-эр-ци-тив-ный
по-лу-не-пре-рыв-ный
ми-ни-ми-зи-ру-ю-щу-ю
по-сле-до-ва-тель-ность
ко-эр-ци-тив-ность
огра-ни-чен-ной
ре-флек-сив-ность
под-по-сле-до-ва-тель-ность
по-лу-не-пре-рыв-ность
иден-ти-фи-ци-ру-е-мо-сти
ин-тер-пре-та-ци-он-ной
функ-ци-о-наль-ный
вос-про-из-во-ди-му-ю
од-но-вре-мен-ном
ин-ва-ри-ан-тен
}
\maketitle

\begin{center}
\fcolorbox{obsblue}{lightblue}{%
\parbox{0.91\linewidth}{\small
\textbf{Статус.} Эталонный документ фазы P0 для системы GO-QS v0.1.
Текст определяет математический интерфейс задачи наблюдения и оценивания.
Он не вводит новую физическую теорию, не объявляет эвристические фильтры
инвариантами природы и не отождествляет результат оптимизации с физической
геометрией без отдельно заданной интерпретационной карты.}}
\end{center}

\begin{abstract}
Задаётся типизированный функциональный интерфейс геометрии наблюдения.
Симметрия или смена кадра, редукция разрешения,
де\-тер\-ми\-ни\-ро\-ван\-ное наблюдение,
стохастический канал и оцениватель разделены как операторы разных типов.
``Нулевая сцена'' заменена явной калибровочной системой; геометрическое
сечение определено через регулярное множество уровня и pullback на вложенную
подповерхность; хроматическая уникальность ограничена конечным конфликтным
графом. Все слагаемые вариационного функционала безразмерны, снабжены
протоколом и политикой нулевого знаменателя. Доказана теорема существования
минимизатора прямым методом и сформулирован критерий структурной
восстановимости. Причинность введена через лоренцево отношение $J^-$, а
Landauer--Bekenstein--Margolus--Levitin ограничения вынесены из универсального
функционала в условные физические модули. Результат является интерфейсом для
воспроизводимых inverse problems, а не самостоятельной динамической теорией.
\end{abstract}

\tableofcontents

\section{Область действия и статусы утверждений}

В документе используются шесть статусов:
\[
\status{DEFINITION},\quad
\status{THEOREM},\quad
\status{PROPOSITION},\quad
\status{MODEL},\quad
\status{HYPOTHESIS},\quad
\status{DIAGNOSTIC}.
\]
Определения и теоремы относятся к математическому ядру. Модель фиксирует
конкретный статистический или вычислительный выбор. Гипотеза требует
независимого обоснования. Диагностика является вычислимым признаком и не
считается физическим инвариантом.

\begin{warning}[Граница новизны]
Карты наблюдения, марковские ядра, регуляризация и прямой метод
вариационного исчисления являются стандартными конструкциями. Содержание
этого текста --- их строгая типовая сборка для корпуса Geometry of
Observation и устранение несовместимых размерностей прежнего интерфейса.
\end{warning}

\section{Нормативный пакет наблюдения GO-QS}

\begin{definition}[Пакет наблюдения]
Пусть $\Lambda$ --- множество допустимых протоколов. Пакетом наблюдения
называется
\[
\mathfrak O=
\left(
X,\Phi,\xi,\Lambda,
\set{X_\lambda,Y_\lambda^0,Y_\lambda,
R_\lambda,O_\lambda,K_\lambda,\widehat\theta_\lambda}_{\lambda\in\Lambda},
G,\Qreg,\Dprof
\right),
\]
где:
\begin{enumerate}[leftmargin=*]
\item $X$ --- скрытое пространство состояний с объявленной категорией;
\item $\Phi^t:X\to X$ --- полная динамика, если она задана;
\item $\xi$ --- дескриптор собственных масштабов системы;
\item $R_\lambda:X\to X_\lambda$ --- редукция или огрубление;
\item $O_\lambda:X_\lambda\to Y_\lambda^0$ --- идеальная
детерминированная карта наблюдения;
\item $K_\lambda(\dd y\mid y^0)$ --- марковское ядро из $Y_\lambda^0$ в
измеримое пространство данных $Y_\lambda$;
\item $\widehat\theta_\lambda:Y_\lambda\rightrightarrows\Theta_\lambda$ ---
возможно многозначный оцениватель;
\item $G$ действует на состояниях обратимыми отображениями $T_g:X\to X$;
\item $\Qreg$ --- реестр типизированных величин;
\item $\Dprof$ --- профиль разнотипных дефектов.
\end{enumerate}
\end{definition}

Условный закон данных при скрытом состоянии $x$ равен
\[
P_\lambda(\dd y\mid x)
=K_\lambda\!\left(\dd y\mid
O_\lambda(R_\lambda x)\right).
\]
Детерминированный датчик получается при ядре Дирака.

\begin{figure}[htbp]
\centering
\begin{tikzpicture}[
  node distance=8mm and 9mm,
  box/.style={draw,rounded corners,minimum height=9mm,align=center,fill=lightblue},
  arr/.style={-{Latex[length=2.2mm]},thick}
]
\node[box] (x) {$X$\\скрытое};
\node[box,right=of x] (xl) {$X_\lambda$\\редуцированное};
\node[box,right=of xl] (y0) {$Y_\lambda^0$\\идеальный выход};
\node[box,right=of y0] (y) {$Y_\lambda$\\данные};
\node[box,below=of y] (th) {$\widehat\theta_\lambda(y)$\\оценка};
\draw[arr] (x) -- node[above] {$R_\lambda$} (xl);
\draw[arr] (xl) -- node[above] {$O_\lambda$} (y0);
\draw[arr] (y0) -- node[above] {$K_\lambda$} (y);
\draw[arr] (y) -- node[right] {estimator} (th);
\draw[arr,obsred,bend left=37] (x) to node[below] {$T_g$ обратимо} (x);
\end{tikzpicture}
\caption{Типы операторов не взаимозаменяемы. $T_g$ не является каналом
потери информации.}
\end{figure}

\begin{proposition}[Композиционный тип]
Если $R_\lambda$ и $O_\lambda$ измеримы, а $K_\lambda$ --- марковское
ядро, то $P_\lambda(\cdot\mid x)$ задаёт марковское ядро из $X$ в
$Y_\lambda$.
\end{proposition}
\begin{proof}
Для каждого $x$ это вероятностная мера как сечение $K_\lambda$.
Измеримость $x\mapsto P_\lambda(A\mid x)$ для измеримого $A$ следует из
измеримости $O_\lambda\circ R_\lambda$ и определяющего свойства ядра.
\end{proof}

\subsection{Паспорт величины}

Каждая величина $q$ имеет тип
\[
\Type(q)=
(D_q,V_q,\kappa_q,\bm d_q,U_q,C_q,\rho_q,\lambda_q,\xi_q,
\mathsf{Unc}_q,\mathsf{Status}_q),
\]
где $\bm d_q\in\mathbb Q^7$ --- SI-вектор размерности, $\kappa_q$ ---
семантический тип, $C_q$ --- контекст единиц, $\rho_q$ --- закон действия
кадра или калибровочной группы. Физически безразмерные угол, вероятность,
ранг, число элементов и информация не считаются автоматически
арифметически совместимыми.

\section{Калибровочная сцена и камера-трасса}

\begin{definition}[Калибровочная сцена]
Калибровочная сцена есть пакет
\[
\mathcal C_0=(o,E_0,\ell_0,t_0,\mathcal W_0),
\]
где $o$ --- начало, $E_0=(e_1,\ldots,e_n)$ --- ориентированный
ортонормированный репер, $\ell_0>0$ и $t_0>0$ --- эталонные длина и время,
а $\mathcal W_0$ --- область валидности калибровки.
\end{definition}

Запись $0!=1$ из v0.1 не является частью определения. Её можно использовать
только как мнемонику ``одна выбранная исходная калибровка''. Ни единственность
репера, ни его физическая привилегированность из этой записи не следуют.

\begin{definition}[Камера-трасса]
При евклидовой пространственной кинематике камера-трасса есть измеримая или
гладкая кривая
\[
C_\lambda:I_{\rm cam}\to SE(n).
\]
Она входит в протокол $\lambda$. Если $C_\lambda$ неизвестна, её параметры
$c$ включаются в пространство $\Theta_\lambda$, а наблюдаемость $c$ должна
проверяться отдельно.
\end{definition}

Смена координат, активное перемещение камеры и физическая симметрия объекта
являются разными операциями. Их отождествление запрещено реестром GO-QS.

\subsection{Строгое содержание ghost-mode}

\begin{definition}[Детерминированный слой совместимости]
Для идеального выхода $y^0\in Y_\lambda^0$ определяется
\[
F_{y^0}^{\lambda}
=\set{x\in X:O_\lambda(R_\lambda x)=y^0}.
\]
\end{definition}

При заданном prior $\pi(\dd x)$ и стохастическом канале скрытая
неопределённость описывается posterior $\pi_\lambda(\dd x\mid y)$.
Термин ``ghost-mode'' далее означает либо слой $F_{y^0}^{\lambda}$, либо
posterior-распределение, но не новый неформальный кортеж невидимых объектов.

\section{Геометрические сечения: корректный тип}

Пусть $M$ --- гладкое $d$-мерное многообразие, а $z_t$ --- поле или
геометрическая структура на $M$.

\begin{definition}[Регулярное сечение уровня]
Для гладкой функции $f_{j,\lambda}:M\to\R$ и её регулярного значения
$c_{j,\lambda}$ положим
\[
\Sigma_{j,\lambda}=f_{j,\lambda}^{-1}(c_{j,\lambda}),\qquad
\iota_{j,\lambda}:\Sigma_{j,\lambda}\hookrightarrow M.
\]
Тогда $\Sigma_{j,\lambda}$ является вложенным подмногообразием размерности
$d-1$, а $\iota_{j,\lambda}$ --- включение.
\end{definition}

Двумерное сечение получается при $d=3$. Для поля $z_t$ корректная операция
ограничения есть pullback $\iota_{j,\lambda}^{*}z_t$ тогда, когда тип поля
допускает pullback. Дискретный выход задаётся отдельно:
\[
S_{j,\lambda}
=Q_{j,\lambda}\circ\iota_{j,\lambda}^{*}:
\mathcal Z(M)\longrightarrow\mathcal Z(Q_{j,\lambda}),
\]
где $Q_{j,\lambda}$ --- конечная сетка, квантователь или оператор выборки.

\begin{remark}
Карта $M\to\Sigma$ сама по себе не называется сечением. Ретракция
$r:M\to\Sigma$, сечение расслоения $s:B\to E$ и включение множества уровня
$\iota:\Sigma\hookrightarrow M$ имеют разные типы.
\end{remark}

Для графа $G=(V,E)$ аналогом является объявленный оператор выбора
$V_{j,\lambda}\subseteq V$ и индуцированный подграф. Он не смешивается с
гладким pullback.

\section{Конечная хроматическая диагностика}

\begin{definition}[Конфликтный граф]
После дискретизации пусть $V_{j,\lambda}$ конечно. Объявленная симметричная
relation $\mathcal R_{\chi,\lambda}$ задаёт
\[
G_{\chi,\lambda}=(V_{j,\lambda},E_{\chi,\lambda}),\qquad
(v,w)\in E_{\chi,\lambda}\iff
v\ne w\ \text{и}\ \mathcal R_{\chi,\lambda}(v,w).
\]
\end{definition}

Для меток $\chi:V_{j,\lambda}\to\mathcal C$ нормированный collision defect
равен
\[
\delta_\chi=
\begin{cases}
\displaystyle
\frac{1}{|E_{\chi,\lambda}|}
\sum_{(v,w)\in E_{\chi,\lambda}}
\mathbf 1_{\{\chi(v)=\chi(w)\}},&
|E_{\chi,\lambda}|>0,\\[3mm]
0,&|E_{\chi,\lambda}|=0.
\end{cases}
\]
Это число из $[0,1]$ и имеет статус \status{DIAGNOSTIC}. Для непрерывных
цветовых признаков разрешена margin-loss только после объявления метрики
$d_{\mathcal C}$ и масштаба $m_{\mathcal C}>0$.

\begin{warning}[Континуальная патология v0.1]
Требование различать все точки с рациональной разностью на континууме не
принимается как базовый постулат. Рациональные разности плотны, а
непрерывный RGB-код не обеспечивает заявленную глобальную уникальность.
Допустимы конечная сетка, конечная $\varepsilon_x$-сеть или отдельная
теорема о раскраске конкретного графа.
\end{warning}

\section{Операциональные масштабы и геометрические фильтры}

\subsection{Масштаб}

Для положительной длины $\ell$ используются
\[
\widehat\ell=\frac{\ell}{\ell_{\rm ref}},
\qquad
\chi_L=\log_b\frac{\ell}{\ell_{\rm ref}},
\quad b>0,\ b\ne1.
\]
Аргумент логарифма безразмерен. Спиральная параметризация
$r(\theta)=r_0e^{a\theta}$ допустима при безразмерном $a\theta$.
Золотое сечение, рациональная сетка или спектральный масштаб являются
опциональными model priors, а не универсальными линейками.

\subsection{Пифагоров диагностический штраф}

Пусть $u,v\ne0$ --- рёбра ячейки и
\[
r_u=\frac{\norm{u}}{\sqrt{\norm{u}^2+\norm{v}^2}},
\qquad
r_v=\frac{\norm{v}}{\sqrt{\norm{u}^2+\norm{v}^2}}.
\]
Для конечного набора примитивных троек
\[
\mathcal P_A=\set{(a,b,c)\in\N^3:
a^2+b^2=c^2,\ \gcd(a,b,c)=1,\ c\le A}
\]
определим
\[
\delta_{\rm pyth}(u,v)=
\min_{(a,b,c)\in\mathcal P_A}
\left[
\left(r_u-\frac ac\right)^2+
\left(r_v-\frac bc\right)^2+
\beta_{\perp}
\left(\frac{\langle u,v\rangle}{\norm{u}\norm{v}}\right)^2
\right].
\]
Если порядок рёбер не фиксирован, минимум берётся также по перестановке
$a,b$. Все слагаемые безразмерны. Нулевые рёбра исключены из допустимой
области.

\subsection{Герон и Кэли--Менгер}

Формулы Герона и Кэли--Менгера являются тождествами евклидовой геометрии,
а не самостоятельными физическими ограничениями. Если они используются
как тест согласованности шумных данных, остаток нормируется:
\[
\delta_{\rm Heron}
=\frac{\abs{A_{\rm obs}^2-s(s-a)(s-b)(s-c)}}{\ell_{\rm ref}^{\,4}},
\qquad
s=\frac{a+b+c}{2}.
\]
Если вводится предпочтение целочисленных или рациональных площадей, оно
помечается \status{MODEL}, задаётся конечным target set в переменной
$A/\ell_{\rm ref}^2$ и исключается из функционала по умолчанию.

\subsection{Симметрия}

Для нормированного представления $z$ в гильбертовом пространстве и
объявленной группы $G_\lambda$ допустима диагностика
\[
\delta_{\rm sym}(z)=
\begin{cases}
\displaystyle
\inf_{g\in G_\lambda\setminus\{e\}}
\frac{\norm{z-g\!\cdot\! z}^{2}}{\norm z^{2}},
&\norm z>0,\ G_\lambda\ne\{e\},\\[2mm]
0,&\norm z=0\ \text{или}\ G_\lambda=\{e\}.
\end{cases}
\]
Группа, представление и нулевая политика являются частью протокола.

\section{Динамический дефект замыкания}

Пусть $y_v:[0,T_{\rm hor}]\to(Y_\lambda,d_{Y,\lambda})$ --- наблюдаемая
траектория маркера, а $r_Y(\lambda)>0$ --- разрешение выхода. Для
$0<P_{\min}\le P\le P_{\max}<T_{\rm hor}$ положим
\[
\delta_{{\rm cl},v}(P)=
\left[
\frac{1}{T_{\rm hor}-P}
\int_0^{T_{\rm hor}-P}
\left(
\frac{d_{Y,\lambda}(y_v(t+P),y_v(t))}
{r_Y(\lambda)}
\right)^2\dd t
\right]^{1/2},
\]
\[
\delta_{{\rm cl},v}
=\inf_{P\in[P_{\min},P_{\max}]}\delta_{{\rm cl},v}(P).
\]
Эта величина безразмерна и протокол-зависима. Малое значение на конечном
горизонте не доказывает периодичность полной скрытой динамики.

\section{Причинность и физические бюджеты}

\subsection{Лоренцева причинность}

Пусть $(M,g)$ --- time-oriented Lorentzian manifold. Данные, считанные в
событии $p_{\rm obs}$, могут зависеть от события $q$ только если
\[
q\in J^-(p_{\rm obs}).
\]
Это hard constraint допустимого множества $\A_\lambda$. В плоской
пространственно-временной диаграмме с координатами $(t,\bm x)$ условие
имеет вид
\[
t_q\le t_{\rm obs},\qquad
\norm{\bm x_{\rm obs}-\bm x_q}
\le c(t_{\rm obs}-t_q).
\]
Риманов шар $d(x,o)\le c|t-t_0|$ не переносится в общий лоренцев случай.

\subsection{Условные информационно-физические модули}

Следующие формулы не входят в универсальный функционал:
\begin{align*}
Q_{\rm erase}&\ge k_BT\ln2\;N_{\rm erased},\\
N_{\rm bits}&\le
\frac{2\pi ER}{\hbar c\ln2},\\
\tau_\perp&\ge
\frac{\pi\hbar}{2(E-E_0)}.
\end{align*}
Первая относится к логически необратимому стиранию при объявленной
термодинамической модели; вторая требует условий применимости
Bekenstein bound; третья --- к времени перехода в ортогональное
квантовое состояние при конечной средней энергии над основным состоянием.
Биты, джоули и секунды не суммируются. При необходимости каждая формула
задаёт отдельное допустимое множество или отдельный нормированный defect.

\subsection{Нет универсальной неопределённости ``сечение--путь''}

Не существует выведенной в v0.1 теоремы
$\Delta_\Sigma\Delta_\gamma\ge\kappa_{\rm obs}$. Для конкретного прибора
разрешена только модель
\[
\widehat u_\Sigma\widehat u_\gamma\ge\kappa_\lambda,\qquad
\widehat u_\Sigma=\frac{u_\Sigma}{u_{\Sigma,0}},\quad
\widehat u_\gamma=\frac{u_\gamma}{u_{\gamma,0}},
\]
где обе нормировки, метод оценки неопределённостей и область валидности
получены из калибровки. Статус такого отношения --- \status{HYPOTHESIS}
или \status{EMPIRICAL}, но не квантовая теорема.

\section{Безразмерный вариационный оцениватель}

\subsection{Единое пространство параметров}

Пусть $\Theta_\lambda$ --- объявленное пространство параметров, а
$\A_\lambda\subseteq\Theta_\lambda$ --- допустимое множество. Один и тот
же символ $\theta\in\A_\lambda$ используется во всех членах. Например,
\[
\theta=(z_0,\vartheta_\Phi,c,\vartheta_R,\vartheta_O,\chi),
\]
но карта наблюдения или камера включается в $\theta$ только если она
действительно оценивается из данных.

\subsection{Согласование с данными}

Для гауссовой модели с положительно определённой ковариацией
$\Sigma_y(t)$:
\[
\widehat D_y(\theta;y)=
\int
\norm{
\Sigma_y(t)^{-1/2}
\bigl(y_\theta(t)-y(t)\bigr)
}^{2}\,\dd\widehat\mu_T(t),
\qquad
\dd\widehat\mu_T=\frac{\dd t}{T_{\rm hor}}.
\]
Для негауссового канала используется безразмерная отрицательная
лог-правдоподобность с явно указанной базовой мерой.

\subsection{Регуляризаторы}

Пусть $S_Z$ --- обратимый scale operator для компонент $z$, а $T_0>0$.
Тогда
\[
\widehat E_1(z)=
\int_0^1
\norm{S_Z^{-1}T_0\,\partial_tz(T_0\widehat t)}^2
\dd\widehat t,
\]
\[
\widehat E_2(z)=
\int_0^1
\norm{S_Z^{-1}T_0^2\,\partial_t^2z(T_0\widehat t)}^2
\dd\widehat t.
\]
Для регулярной кривой длины $L_{\rm path}>0$:
\[
\widehat E_\kappa=
\frac{1}{L_{\rm path}}\int
(\ell_0\kappa(s))^2\,\dd s.
\]
Все три величины безразмерны.

\subsection{Скаляризация}

Пусть $\mathcal K_\lambda$ --- конечный набор активных диагностик,
$\delta_k=\nu_k(D_k)\in[0,\infty]$ --- объявленные нормировки, а
$\phi_k:[0,\infty]\to[0,\infty]$ --- объявленные penalty maps. Определим
\[
\boxed{
J_\lambda(\theta;y)=
w_y\widehat D_y
w_1\widehat E_1
w_2\widehat E_2
w_\kappa\widehat E_\kappa
\sum_{k\in\mathcal K_\lambda}w_k\phi_k(\delta_k)
I_{\A_\lambda}(\theta)
}
\]
с $w_\bullet\ge0$ и суммой весов $1$. Здесь $I_{\A_\lambda}$ равно
$0$ на $\A_\lambda$ и $+\infty$ вне неё. Скаляризация является
decision functional, а не универсальным геометрическим инвариантом.

\begin{proposition}[Инвариантность при смене единиц]
Если каждый член $J_\lambda$ построен из безразмерных отношений,
нормированных мер и согласованно преобразованных ковариаций, то значение
$J_\lambda$ не меняется при согласованной замене единиц.
\end{proposition}
\begin{proof}
Каждое отношение физической величины к эталону той же размерности
инвариантно. Mahalanobis residual инвариантен при одновременном
преобразовании данных и ковариации. Нормированные меры безразмерны.
Следовательно, инвариантна каждая сумма и вся их безразмерная линейная
комбинация.
\end{proof}

\begin{definition}[Вариационный оцениватель]
\[
\widehat\theta_\lambda(y)
=\Argmin_{\theta\in\Theta_\lambda}J_\lambda(\theta;y).
\]
Значение является множеством, пока единственность не доказана.
\end{definition}

\section{Существование, неединственность и интерпретация}

\begin{theorem}[Прямой метод]
Пусть $V$ --- рефлексивное банахово пространство,
$\A_\lambda\subset V$ непусто и последовательно слабо замкнуто. Пусть
$J_\lambda:V\to(-\infty,+\infty]$ собственный, ограниченный снизу,
коэрцитивный и последовательно слабо нижне полунепрерывный. Тогда
\[
\Argmin_{V}J_\lambda\ne\varnothing.
\]
\end{theorem}
\begin{proof}
Возьмём минимизирующую последовательность. Коэрцитивность делает её
ограниченной. Рефлексивность даёт слабо сходящуюся подпоследовательность.
Слабая замкнутость сохраняет предел в $\A_\lambda$, а слабая нижняя
полунепрерывность даёт значение не больше $\liminf$ значений
последовательности. Предел достигает инфимума.
\end{proof}

\begin{warning}
Теорема не даёт единственности, устойчивости, идентифицируемости или
физической истинности forward model. Для каждого конкретного пространства
необходимо отдельно доказать коэрцитивность и слабую нижнюю
полунепрерывность.
\end{warning}

\begin{definition}[Интерпретационная карта]
Реконструированная геометрическая структура задаётся объявленной картой
\[
\mathcal G_\lambda:\Theta_\lambda\to\mathsf{Geom}_\lambda.
\]
Выходом задачи является множество
$\mathcal G_\lambda(\widehat\theta_\lambda(y))$, а не автоматически
``геометрия пространства''.
\end{definition}

\section{Восстановимость}

\begin{theorem}[Факторизация скалярной величины]
Пусть $O:X\to Y$ --- отображение множеств и $q:X\to V$. Существует
$\overline q:O(X)\to V$ такое, что
\[
q=\overline q\circ O
\]
тогда и только тогда, когда $q$ постоянна на каждом слое $O^{-1}(y)$.
Карта $\overline q$ единственна на $O(X)$.
\end{theorem}
\begin{proof}
Необходимость следует из равенства $O(x)=O(x')$. Для достаточности
положим $\overline q(y)=q(x)$ для любого $x\in O^{-1}(y)$.
Постоянство на слоях обеспечивает корректность определения.
\end{proof}

В стохастическом случае условие $H_b(q(X)\mid Y)=0$ означает
восстановимость почти наверное только относительно объявленного
распределения и не заменяет структурную теорему.

\section{Quantity and Protocol Ledger}

\subsection{Обязательный протокол}

\begin{longtable}{L{0.27\linewidth}L{0.64\linewidth}}
\toprule
Поле & Нормативное содержание\\
\midrule
\endhead
\texttt{hidden\_space} & категория $X$, топология или $\sigma$-алгебра,
область динамики\\
\texttt{system\_descriptor} & $L_X,T_X,M_X,E_X$, число степеней свободы
и прочие собственные масштабы\\
\texttt{reduction} & сигнатура $R_\lambda:X\to X_\lambda$ и параметры
разрешения\\
\texttt{observation} & сигнатура $O_\lambda:X_\lambda\to Y_\lambda^0$\\
\texttt{noise\_kernel} & $K_\lambda(\dd y\mid y^0)$, калибровка и
ковариация\\
\texttt{frame} & координаты, базис, camera pose и закон преобразования\\
\texttt{quantizer} & сетка, ADC, partition или правило missing data\\
\texttt{horizon} & $T_{\rm hor}$ и temporal resolution $\delta t$\\
\texttt{estimator} & $\Theta_\lambda$, loss, regularizers, веса и solver\\
\texttt{uncertainty} & covariance, interval, posterior или risk bound\\
\texttt{defects} & кодомен каждого $D_k$ и нормировка $\nu_k$\\
\bottomrule
\end{longtable}

\subsection{Ключевые величины}

\begin{longtable}{L{0.18\linewidth}L{0.19\linewidth}L{0.18\linewidth}L{0.34\linewidth}}
\toprule
Величина & Размерность $\bm d$ & Тип & Условие\\
\midrule
\endhead
$\ell_0,\ell_{\rm ref},\varepsilon_x$ & $(1,0,0,0,0,0,0)$ &
scalar & положительны\\
$T_0,T_{\rm hor},\delta t$ & $(0,0,1,0,0,0,0)$ &
scalar & положительны\\
$\kappa$ & $(-1,0,0,0,0,0,0)$ &
scalar & требует smoothing protocol\\
$\Sigma_y$ & квадрат единицы $y$ &
covariance tensor & положительно определена или regularized\\
$\delta_k$ & $\bm 0$ & loss & нормировка и zero policy обязательны\\
$J_\lambda$ & $\bm 0$ & loss & decision functional\\
$H_b$ & $\bm 0$ & information & основание $b$ обязательно\\
$N_{\rm bits}$ & $\bm 0$ & information &
отображение count $\to$ bit объявляется\\
\bottomrule
\end{longtable}

\section{Журнал исправлений: v0.1 к v0.2}

\begin{enumerate}[leftmargin=*]
\item Generic $\mathcal S_j:M\to\Sigma_j$ заменён регулярным множеством
уровня, включением $\iota_j$ и дискретизатором $Q_j$.
\item Пакет параметров унифицирован как $\theta\in\Theta_\lambda$.
\item Камера перенесена в протокол; неизвестная камера допускается только
как оцениваемый параметр.
\item Ghost-mode заменён слоем совместимости или posterior.
\item Континуальная рациональная раскраска заменена конечным конфликтным
графом.
\item Пифагоров штраф полностью безразмерен; ортогональный остаток
нормирован.
\item Все логарифмы содержат отношение к эталону.
\item Closure defect нормирован разрешением выхода и горизонтом.
\item Риманов ``причинный шар'' заменён отношением $J^-$.
\item Информационные пределы не суммируются с энергиями и геометрическими
остатками.
\item Невыведенная uncertainty product удалена из теоремного ядра.
\item Функционал состоит только из безразмерных членов.
\item $\Argmin$ оставлен множеством; геометрия получается отдельной
интерпретационной картой.
\item Компактностная эвристика заменена точной теоремой прямого метода.
\end{enumerate}

\section{Claim firewall и критерий эталонности P0}

\textbf{Доказано или определено:}
\begin{enumerate}[leftmargin=*]
\item типизированная композиция редукции, observation map и noise kernel;
\item корректный тип геометрического сечения;
\item безразмерность всех активных членов $J_\lambda$;
\item существование minimizer при гипотезах прямого метода;
\item критерий структурной факторизации величины;
\item протокол-зависимый характер диагностик.
\end{enumerate}

\textbf{Не доказано и не заявляется:}
\begin{enumerate}[leftmargin=*]
\item единственность или устойчивость любой конкретной реконструкции;
\item универсальность Pythagorean, Heronian, chromatic или
golden-ratio priors;
\item универсальная uncertainty relation;
\item вывод квантовой механики, гравитации или термодинамики;
\item физическая реальность любой структуры только из факта минимизации.
\end{enumerate}

Документ проходит P0, если машинный manifest подтверждает уникальность
символов в каждой области, полные сигнатуры карт, статусы утверждений и
отсутствие сложения сырых разнотипных дефектов.

\section{Заключение}

Функциональный интерфейс теперь имеет строгую роль: он задаёт
воспроизводимую задачу наблюдения и оценивания
\[
x
\xrightarrow{R_\lambda}
x_\lambda
\xrightarrow{O_\lambda}
y_\lambda^0
\xrightarrow{K_\lambda}
y
\xrightarrow{\widehat\theta_\lambda}
\widehat\theta,
\]
а затем отделяет результат оценивания от его геометрической интерпретации.
Это минимально необходимая архитектура, на которую можно переносить
информационный, метрический, тензорный и прикладные модули без смешения
единиц и математических типов.

{\small
\begin{thebibliography}{99}
\bibitem{LeeSmooth}
J.~M. Lee,
\emph{Introduction to Smooth Manifolds}, 2nd ed.,
Springer, 2013.

\bibitem{Bogachev}
V.~I. Bogachev,
\emph{Measure Theory}, Vols. I--II,
Springer, 2007.

\bibitem{EkelandTemam}
I.~Ekeland and R.~Temam,
\emph{Convex Analysis and Variational Problems},
SIAM, 1999.

\bibitem{Engl}
H.~W. Engl, M.~Hanke, and A.~Neubauer,
\emph{Regularization of Inverse Problems},
Kluwer, 1996.

\bibitem{KaipioSomersalo}
J.~Kaipio and E.~Somersalo,
\emph{Statistical and Computational Inverse Problems},
Springer, 2005.

\bibitem{ONeill}
B.~O'Neill,
\emph{Semi-Riemannian Geometry with Applications to Relativity},
Academic Press, 1983.

\bibitem{Landauer}
R.~Landauer,
Irreversibility and heat generation in the computing process,
\emph{IBM Journal of Research and Development} \textbf{5} (1961),
183--191.

\bibitem{Bekenstein}
J.~D. Bekenstein,
Universal upper bound on the entropy-to-energy ratio for bounded systems,
\emph{Physical Review D} \textbf{23} (1981), 287--298.

\bibitem{MargolusLevitin}
N.~Margolus and L.~B. Levitin,
The maximum speed of dynamical evolution,
\emph{Physica D} \textbf{120} (1998), 188--195.
\end{thebibliography}
}

\end{document}
